\documentclass[11pt]{article}
\usepackage{amsmath, amssymb}
\usepackage[margin=1in]{geometry}

\newcommand{\vv}[1]{\vec{#1}}
\newcommand{\half}{\tfrac{1}{2}}
\DeclareMathOperator{\atantwo}{atan2}

\title{Rounded polygons: overlap area (6a / 6b)}
\date{}
\author{}

\begin{document}
\maketitle

Polygons $A,B$ intersect and generate a region $A\cap B$ with boundary $\partial(A\cap B)$. Its
area is the Green's-theorem boundary integral
\[
  |A\cap B| = \half\,\varepsilon_{\alpha\beta}\oint_{\partial(A\cap B)} X^{\alpha}\,dX^{\beta},
\]
accumulated feature by feature relative to a reference point $\vv R$ = the first backbone vertex
of polygon $A$ (the lower index, $A<B$). The total is $\vv R$-invariant -- the $\vv R$ terms
telescope around the closed loop -- so $\vv R$ is only there to keep the arithmetic local.

\section{Per-feature area}

A feature piece from point $\vv P$ to $\vv Q$ contributes, for an \textbf{edge},
\[
  h_e(\vv P,\vv Q)=\half\,\varepsilon_{\alpha\beta}(P-R)^{\alpha}(Q-R)^{\beta},
\]
and for an \textbf{arc} (center $\vv z$, radius $\rho$) a chord term plus a circular segment,
\[
  h_a(\vv P,\vv Q)=\half\,\varepsilon_{\alpha\beta}(P-R)^{\alpha}(Q-R)^{\beta}
  +\half\rho^{2}(\theta-\sin\theta),\qquad
  \theta=\atantwo\!\Big(\varepsilon_{\alpha\beta}(P-z)^{\alpha}(Q-z)^{\beta},\;(P-z)\!\cdot\!(Q-z)\Big),
\]
where $\theta$ is the signed central angle of the sub-arc $\vv P\to\vv Q$. The feature selector is
\[
  h(\vv P,\vv Q)=\begin{cases} h_e(\vv P,\vv Q) & f\ \text{odd (edge)}\\[2pt]
                               h_a(\vv P,\vv Q) & f\ \text{even (arc).}\end{cases}
\]

\section{Traversal}

The overlap boundary is the kept runs of $\partial A$ inside $B$ and $\partial B$ inside $A$,
each from an intersection to its outersection. Label an intersection pair $P=\{i,j,k,l\}$: feature
$i$ (on $A$) meets feature $j$ (on $B$), and following $i$ CCW reaches feature $l$, which meets
$k$ (shapes $\{A,B,B,A\}$). For an intersection at boundary coordinate $\sigma$, let
$\sigma^{-}\equiv\lfloor\sigma\rfloor$ be its feature index (the kiss point starting that
feature), and let $z(\cdot),z'(\cdot)$ be the next / previous feature in CCW order (so
$\vv q_{z(\sigma^{-})}$ is that feature's far kiss point). With $\vv q$ the intersection and kiss
points,
\[
  U=\sum_{P=\{i,j,k,l\}}\Big\{\,\delta_{i^{-}=l^{-}}\,h(\vv q_i,\vv q_l)
  +(1-\delta_{i^{-}=l^{-}})\Big[\,h(\vv q_i,\vv q_{z(i^{-})})
  +\!\!\sum_{m=z(i^{-})}^{z'(l^{-})}\! h(\vv q_m,\vv q_{z(m)})
  +h(\vv q_{l^{-}},\vv q_l)\,\Big]\Big\}.
\]
If the run sits on one feature, $h(\vv q_i,\vv q_l)$; otherwise the partial start feature
($\vv q_i$ to its end $\vv q_{z(i^{-})}$), the whole interior features ($\vv q_m$ to
$\vv q_{z(m)}$), and the partial end feature ($\vv q_{l^{-}}$ to $\vv q_l$).

\section{Parallel split, $U=U_{\mathrm{out}}+U_{\mathrm{in}}$}

The interior sum is awkward to parallelize, so peel it off:
\[
  U_{\mathrm{out}}=\sum_{P}\Big[\delta_{i^{-}=l^{-}}h(\vv q_i,\vv q_l)
  +(1-\delta_{i^{-}=l^{-}})\big(h(\vv q_i,\vv q_{z(i^{-})})+h(\vv q_{l^{-}},\vv q_l)\big)\Big],\qquad
  U_{\mathrm{in}}=\sum_{m=0}^{2n}\;\sum_{P}\delta_{\,m\,\in\,\{z(i^{-}),\,\dots,\,z'(l^{-})\}}\;h(\vv q_m,\vv q_{z(m)}).
\]
$U_{\mathrm{out}}$ is per intersection pair (the two partial end features); $U_{\mathrm{in}}$ is
per \emph{feature} $m$ (the whole interior features), which distributes over threads. The test
is membership $m\in\{z(i^{-}),\dots,z'(l^{-})\}$ --- the features strictly CCW-between $i$ and
$l$, walked by $z$. This is the right criterion: it respects the cyclic feature order, whereas
an $i^{-}<m<l^{-}$ inequality silently drops the interior whenever a run wraps past feature
$0$.

\paragraph{Validation.} Built both ways; the walk and the $U_{\mathrm{out}}/U_{\mathrm{in}}$
split agree to machine precision, and the area matches a Monte-Carlo estimate to $\sim0.1\%$
(convex, free and periodic) and to $\sim0.5\%$ across $60$ non-convex star-pair overlaps (up to
$8$ intersections).

\end{document}
